%---------------------------------------------------------------------------------------

\chapter{Implementação Computacional}\label{ch:implementacao-computacional}

%---------------------------------------------------------------------------------------

% Lista de abreviaturas e de siglas
\criarsigla{AoS}{\textit{Array of Structures}}
\criarsigla{API}{\textit{Application Programming Interface}}
\criarsigla{CUDA}{\textit{Compute Unified Device Architecture}}
\criarsigla{FLOPs}{\textit{Floating Point Operations}}
\criarsigla{IDW}{\textit{Inverse Distance Weighting}}
\criarsigla{JIT}{\textit{Just-In-Time}}
\criarsigla{NUMA}{\textit{Non-Uniform Memory Access}}
\criarsigla{RAM}{\textit{Random Access Memory}}
\criarsigla{SDK}{\textit{Software Development Kit}}
\criarsigla{SIMD}{\textit{Single Instruction, Multiple Data}}
\criarsigla{SIMT}{\textit{Single Instruction, Multiple Threads}}
\criarsigla{SM}{\textit{Streaming Multiprocessor}}
\criarsigla{SoA}{\textit{Structure of Arrays}}
\criarsigla{VRAM}{\textit{Video Random Access Memory}}
\criarsigla{VTK}{\textit{Visualization Toolkit}}
\criarsigla{YAML}{\textit{Yet Another Markup Language}}

%---------------------------------------------------------------------------------------

% Lista de simbolos

%---------------------------------------------------------------------------------------

A tradução da formulação apresentada no capítulo anterior em um programa executável envolve decisões que afetam o desempenho em arquiteturas heterogêneas, a portabilidade entre fabricantes de \textit{hardware} e a facilidade de manutenção e de extensão do código.
O miniaplicativo de dinâmica dos fluidos computacional (CFD) desenvolvido neste trabalho foi concebido para que essas decisões pudessem ser isoladas, mantendo-se inalterado o método de Lattice Boltzmann (LBM) para escoamentos compressíveis e variando-se apenas a tecnologia responsável por executá-lo.
O mesmo núcleo de cálculo foi reimplementado em cada uma das tecnologias de paralelização consideradas, a saber, OpenMP, OpenACC, CUDA, OpenCL, Kokkos, RAJA e MPI, tomando como referência a versão sequencial.
Para assegurar uma comparação justa entre as implementações, adota-se como critério a invariância da física do modelo, da malha computacional e da disposição dos dados em memória.
Dessa forma, garante-se que as variações de desempenho e de complexidade observadas decorram exclusivamente das características intrínsecas ao modelo de execução e à gestão de recursos de cada tecnologia.

O \textit{software} adota uma arquitetura orientada a objetos em C++ moderno, estruturada para abstrair as particularidades dos diferentes modelos de programação de computação de alto desempenho (HPC).
O núcleo da simulação é definido por uma interface comum que estabelece o contrato para os diversos modelos de paralelização, permitindo que um orquestrador central instancie a estratégia adequada em tempo de execução, com base na configuração fornecida pelo usuário.

Uma classe controladora gerencia o ciclo de vida da simulação, a entrada e a saída de dados e a medição de desempenho.
Essa abstração permite que a lógica de controle, como a marcha no tempo apresentada no ALGORITMO~\ref{alg:marcha-tempo-detalhado} e a reconstrução dos campos macroscópicos para exportação, permaneça independente do \textit{hardware} utilizado, seja ele uma CPU em execução sequencial, uma GPU ou um agrupamento de computadores distribuído.

\begin{algorithm}[H]
    \caption{Marcha no tempo para o LBM compressível com acoplamento térmico}\label{alg:marcha-tempo-detalhado}
    \Entrada{Dados do fluido, malha computacional, tempo final $t_{end}$}
    \Saida{Campos macroscópicos $\rho, \mathbf{u}, T$ em cada passo}
    Inicializar campos $\rho, \mathbf{u}, T$ e distribuições $f_i, g_i$\;
    \Enquanto{$t < t_{end}$}{
        \Para{cada célula da malha}{
            Reconstruir $(\rho;\ \mathbf{u};\ T)$ a partir de $f_i$ e $g_i$\;
            \tcp{Newton-Raphson para o modelo térmico}
            Resolver iterativamente para as variáveis macroscópicas\;
            Calcular equilíbrios $f_i^{eq}(\rho, \mathbf{u}, T)$ e $g_i^{eq}(\rho, \mathbf{u}, T)$\;
            Aplicar colisão (BGK): $f_i^* = f_i - \omega_f (f_i - f_i^{eq})$ e $g_i^* = g_i - \omega_g (g_i - g_i^{eq})$\;
        }
        \If{a estratégia de paralelismo é MPI}{
            Sincronizar as zonas de troca de $f_i^*$ e $g_i^*$ entre processos\;
        }
        \Para{cada célula da malha}{
            Propagar: $f_i(\mathbf{x} + \mathbf{c}_i \Delta t, t + \Delta t) = f_i^*(\mathbf{x}, t)$\;
            $g_i(\mathbf{x} + \mathbf{c}_i \Delta t, t + \Delta t) = g_i^*(\mathbf{x}, t)$\;
            Aplicar condições de contorno\;
        }
        $t \leftarrow t + \Delta t$\;
        \If{passo de gravação}{
            Exportar resultados $(\rho;\ \mathbf{u};\ T)$ em formato VTK\;
        }
    }
\end{algorithm}
\fonte{O autor (2026).}

O processo de compilação, discutido na Seção~\ref{sec:compilacao-configuracao}, detecta automaticamente a presença de compiladores e de bibliotecas como o NVIDIA HPC SDK, empregado para OpenACC e OpenMP, o NVCC, empregado para CUDA, e implementações de MPI e de OpenCL.
Isso garante que o miniaplicativo possa ser compilado e executado em uma ampla variedade de ambientes.
O controle da simulação é feito por um arquivo de configuração em formato YAML (\textit{Yet Another Markup Language}), no qual se definem as propriedades do fluido, a malha computacional e a tecnologia de execução.

%---------------------------------------------------------------------------------------

\section{Arquitetura do Código}\label{sec:arquitetura-simulador}

A arquitetura do código foi projetada para suportar a execução em múltiplas tecnologias de paralelização sem que a lógica física do modelo LBM precise ser reescrita.
Para alcançar esse objetivo, o sistema utiliza o padrão de projeto \textit{Strategy}, centrado na interface abstrata \texttt{ILbmTube} e na classe orquestradora \texttt{Executor}, cuja organização é apresentada na \autoref{fig:diagrama-de-classes}.

\begin{figure}[H]
    \centering
    \caption{Organização das classes do miniaplicativo segundo o padrão de projeto \textit{Strategy}, com a interface \texttt{ILbmTube} e suas implementações por tecnologia de paralelização.}
    \label{fig:diagrama-de-classes}
    \fbox{\begin{minipage}{0.85\textwidth}
        \centering
        \vspace{2cm}
        [Espaço reservado para figura: diagrama de classes simplificado da arquitetura.
        No topo, a classe \texttt{Executor}, associada ao leitor do arquivo de configuração e ao exportador de resultados.
        Abaixo, a interface \texttt{ILbmTube} e, dela derivadas, as implementações \textit{Plain}, OpenMP, OpenACC, CUDA, OpenCL, Kokkos, RAJA e MPI, cada uma com suas variantes de CPU e de GPU.
        Ao lado, as classes auxiliares \texttt{ILattice}, \texttt{LatticeD2Q9} e \texttt{DeviceBuffer}.]
        \vspace{2cm}
    \end{minipage}}
    \fonte{O autor (2026).}
\end{figure}

A interface \texttt{ILbmTube} define o contrato que todos os resolvedores devem implementar, incluindo os métodos de inicialização (\texttt{Init}), de execução de passos de tempo (\texttt{Step} e \texttt{Run}) e de acesso aos campos macroscópicos (\texttt{P}, \texttt{Rho}, \texttt{T} e \texttt{U}).
Essa abstração permite que o \texttt{Executor} gerencie o ciclo de vida da simulação de forma agnóstica ao \textit{hardware}, conforme apresentado no PROGRAMA~\ref{lst:interface-ilbmtube}.

\noindent\begin{minipage}{\linewidth}
\begin{lstlisting}[language=C++, caption={Interface base ILbmTube, que define o contrato dos resolvedores.}, label={lst:interface-ilbmtube}]
template<typename Scalar, typename LatticeType>
class ILbmTube {
public:
  virtual ~ILbmTube() = default;
  virtual Tensor<Scalar, LatticeType::Dim()> P() const = 0;
  virtual Tensor<Scalar, LatticeType::Dim()> Rho() const = 0;
  virtual Tensor<Scalar, LatticeType::Dim()> T() const = 0;
  virtual Tensor<Scalar, LatticeType::Dim() + 1> U() const = 0;
  virtual void Init() = 0;
  virtual void Step(bool save) = 0;
  virtual void Run(Index steps, bool save) = 0;
};
\end{lstlisting}
\fonte{O autor (2026).}
\end{minipage}

O \texttt{Executor} seleciona a implementação de \texttt{ILbmTube} em tempo de execução a partir do arquivo de configuração do qual o usuário tem acesso, sendo também responsável pela medição de desempenho e pela exportação dos resultados em formato VTK para visualização posterior.
O efeito dessa arquitetura sobre a complexidade do código é favorável, pois isola o uso das interfaces de programação de paralelismo em classes específicas, o que facilita a inclusão de novas estratégias.

Componente crítico para o desempenho é a classe \texttt{DeviceBuffer}, responsável pela alocação de memória.
Em arquiteturas de CPU, ela comporta-se como um \texttt{std::vector} alinhado, enquanto em GPUs encapsula ponteiros de memória do dispositivo.
A escolha da disposição de memória por estrutura de vetores (\textit{Structure of Arrays}, SoA) é fundamental: ao armazenar cada componente da distribuição ($f_0, f_1, \dots, f_8$) em blocos contíguos, maximiza-se o aproveitamento do \textit{cache} em CPUs por meio da busca antecipada (\textit{prefetching}), mecanismo pelo qual o processador carrega dados da memória principal para o \textit{cache} antes de sua solicitação explícita.
Essa disposição garante ainda o acesso coalescido à memória (\textit{memory coalescing}) em GPUs, técnica que combina em uma única transação os múltiplos acessos à memória global realizados por diferentes fluxos de execução, desde que os endereços acessados sejam adjacentes.
A \autoref{fig:disposicoes-aos-soa} contrasta essa disposição com a alternativa por vetor de estruturas (\textit{Array of Structures}, AoS).

\begin{figure}[H]
    \centering
    \caption{Disposições de memória AoS e SoA para as funções de distribuição da rede D2Q9 e seu efeito sobre o padrão de acesso simultâneo à memória.}
    \label{fig:disposicoes-aos-soa}
    \fbox{\begin{minipage}{0.85\textwidth}
        \centering
        \vspace{2cm}
        [Espaço reservado para figura: comparação entre as disposições de memória AoS e SoA para as nove distribuições da rede D2Q9.
        Na parte superior, a disposição AoS, com as nove distribuições de uma mesma célula armazenadas de forma contígua.
        Na parte inferior, a disposição SoA, com todas as células de uma mesma direção armazenadas de forma contígua.
        Setas devem indicar os endereços acessados simultaneamente por fluxos de execução vizinhos, evidenciando o acesso coalescido apenas no caso SoA.]
        \vspace{2cm}
    \end{minipage}}
    \fonte{O autor (2026).}
\end{figure}

A marcha no tempo segue uma estrutura de estágios encadeados que minimiza as transferências entre a memória do sistema principal, correspondente à memória RAM, e a memória do dispositivo, correspondente à memória VRAM da GPU.

%---------------------------------------------------------------------------------------

\section{Custo Computacional}\label{sec:custo-computacional}

As decisões de implementação descritas nas seções seguintes somente se justificam frente ao custo do ciclo de cálculo estabelecido na Seção~\ref{subsec:lbm-ciclo-compressivel} do Capítulo~\ref{ch:metodo-de-lattice-boltzmann}, razão pela qual esse custo é quantificado antes de qualquer tecnologia ser apresentada.
Duas grandezas o caracterizam: o número de operações de ponto flutuante (\textit{Floating Point Operations}, FLOPs) executadas por nó e por passo de tempo e o volume de dados movimentado entre a memória e os núcleos de processamento.
A razão entre ambas, denominada intensidade aritmética, determina qual dos dois recursos limita o desempenho e, por consequência, sobre qual deles as otimizações devem incidir.

Na formulação de dupla distribuição sobre a rede D2Q9, cada nó armazena nove valores de $f_i$ e nove de $g_i$.
A atualização de um passo de tempo exige a leitura e a escrita desses dezoito valores, o que corresponde a $288$ \textit{bytes} em precisão dupla, aos quais se somam os campos macroscópicos de massa específica, de velocidade, de temperatura e de pressão.
A propagação interpolada acrescenta a parcela mais expressiva, pois exige a leitura, para cada uma das nove direções, de quatro índices de origem e de quatro pesos de interpolação.
A \autoref{tab:custo-por-no} discrimina essas contribuições para a versão que armazena as tabelas de propagação em memória.

\begin{table}[H]
    \centering
    \caption{Volume de dados movimentado por nó e por passo de tempo em precisão dupla}
    \label{tab:custo-por-no}
    \begin{tabular}{lcc}
        \toprule
        Contribuição & Escalares & \textit{Bytes} \\
        \midrule
        Distribuições $f_i$ e $g_i$ (leitura e escrita) & $36$ & $288$ \\
        Campos macroscópicos $\rho$, $\mathbf{u}$, $T$ e $p$ (leitura e escrita) & $10$ & $80$ \\
        Índices de origem da propagação (leitura)                             & $36$ & $144$ \\
        Pesos de interpolação (leitura)                                       & $36$ & $288$ \\
        \midrule
        Total aproximado                                                      & --- & $800$ \\
        \bottomrule
    \end{tabular}
    \fonte{O autor (2026).}
\end{table}

O tráfego total aproxima-se, portanto, de oitocentos \textit{bytes} por nó e por passo, valor que sequer considera as falhas de \textit{cache} associadas ao acesso não contíguo aos nós de origem.
Os índices de origem e os pesos de interpolação envolvem o mesmo número de escalares por nó, trinta e seis em cada caso, mas contribuem com volumes distintos de dados, respectivamente $144$ e $288$ \textit{bytes}, porque os primeiros são armazenados como inteiros de quatro \textit{bytes} e os segundos em precisão dupla.

O número de operações aritméticas, por sua vez, situa-se na ordem de algumas centenas por nó, uma vez que o cálculo das duas distribuições de equilíbrio e a solução iterativa dos multiplicadores de Lagrange elevam consideravelmente o custo em relação ao LBM isotérmico clássico.
As exponenciais avaliadas nas três iterações de Newton-Raphson, em particular, respondem por uma fração relevante desse total, e seu custo unitário é muito superior ao de uma multiplicação, o que torna o balanço sensível à qualidade da biblioteca matemática empregada por cada compilador.
A leitura desse balanço é usualmente feita pelo modelo \textit{roofline}, que representa, em eixos logarítmicos, o desempenho alcançável em operações por segundo em função da intensidade aritmética do núcleo de cálculo \parencite{Williams2009}.
O modelo estabelece que o desempenho é limitado pelo menor entre dois tetos: o produto da intensidade aritmética pela largura de banda de memória, que aparece como uma reta inclinada, e o desempenho máximo de ponto flutuante do processador, que aparece como um patamar horizontal.
A abscissa em que as duas retas se encontram, denominada ponto de inflexão, separa os núcleos de cálculo limitados pela largura de banda, situados à esquerda, daqueles limitados pela capacidade aritmética, situados à direita.
A \autoref{fig:roofline} representa esse diagrama e a posição estimada do resolvedor sobre ele.

\begin{figure}[H]
    \centering
    \caption{Modelo \textit{roofline} e posição estimada do núcleo de cálculo do miniaplicativo.}
    \label{fig:roofline}
    \fbox{\begin{minipage}{0.85\textwidth}
        \centering
        \vspace{2cm}
        [Espaço reservado para figura: diagrama do modelo \textit{roofline} em eixos logarítmicos, com a intensidade aritmética em operações por \textit{byte} nas abscissas e o desempenho em operações de ponto flutuante por segundo nas ordenadas.
        A reta inclinada, cuja declividade corresponde à largura de banda de memória, deve encontrar o patamar horizontal do desempenho máximo de ponto flutuante no ponto de inflexão, situado entre cinco e dez operações por \textit{byte}.
        Um marcador deve indicar a posição do núcleo de cálculo do miniaplicativo, abaixo de duas operações por \textit{byte}, portanto na região limitada pela largura de banda, com a região à direita do ponto de inflexão identificada como limitada pela capacidade aritmética.]
        \vspace{2cm}
    \end{minipage}}
    \fonte{O autor (2026), com base em \textcite{Williams2009}.}
\end{figure}

A intensidade aritmética resultante é superior à do LBM clássico, mas permanece abaixo de duas operações por \textit{byte}, ao passo que o ponto de inflexão das arquiteturas atuais situa-se tipicamente entre cinco e dez operações por \textit{byte}.
O resolvedor encontra-se, portanto, na região inclinada do diagrama, o que significa que é limitado pela largura de banda de memória, e não pela capacidade aritmética do processador.

Essa constatação orienta três decisões que se repetem nas implementações específicas para cada tecnologia de paralelização.
A primeira é a disposição de memória por estrutura de vetores, que assegura o aproveitamento integral de cada linha de \textit{cache} em CPUs e o acesso coalescido em GPUs, conforme discutido na Seção~\ref{sec:arquitetura-simulador}.
A segunda é a fusão de etapas em um mesmo percurso da malha, que mantém resultados intermediários em registradores em lugar de escrevê-los e reler-los na memória principal.
A terceira é a substituição, nas versões voltadas a aceleradores, das tabelas de propagação pelo cálculo dos índices e dos pesos em tempo de execução.
Essa última decisão é particularmente ilustrativa do compromisso em jogo, pois troca deliberadamente operações aritméticas adicionais pela eliminação de $432$ dos oitocentos \textit{bytes} da \autoref{tab:custo-por-no}, o que é vantajoso precisamente porque o recurso escasso é a largura de banda, e não a capacidade de cálculo.
A verificação experimental dessas expectativas, por meio das taxas de atualização de nós medidas em cada tecnologia, é apresentada no Capítulo~\ref{ch:resultados}.

%---------------------------------------------------------------------------------------

\section{Implementação Sequencial}\label{sec:implementacao-computacional-sequencial}

A versão sequencial, designada no código como \textit{Plain}, estabelece a referência funcional do projeto.
Ela implementa o modelo LBM compressível segundo a formulação de \textcite{Tran2022}, utilizando a classe \texttt{DeviceBuffer} para a alocação de memória e tensores da biblioteca Eigen para a exportação de dados.
Essa implementação fornece a validação científica necessária e serve como modelo de arquitetura para as variantes paralelas.
A estrutura de dados e o fluxo de execução são projetados para facilitar a transposição para aceleradores, evitando construções intrínsecas de CPU que dificultariam a portabilidade.

O modelo de execução é estritamente sequencial, processando cada célula da malha de forma iterativa em um único fluxo de execução.
A gestão de memória ocorre integralmente no sistema principal (RAM), utilizando alinhamento de 64 \textit{bytes} para favorecer a vetorização automática pelo compilador.
Adota-se a disposição SoA, na qual as funções de distribuição $f_i$ e $g_i$ para cada uma das $Q=9$ velocidades da rede D2Q9 são armazenadas em blocos contíguos de tamanho $N$.
Essa escolha maximiza a localidade de dados e permite o uso eficiente das linhas de \textit{cache}.

O mapeamento do algoritmo LBM divide-se em etapas principais: cálculo das grandezas macroscópicas, colisão e propagação.
Na colisão, determinam-se as distribuições de equilíbrio $f_i^{eq}$ e $g_i^{eq}$, aplicando-se o operador BGK.
Para o equilíbrio térmico $g_i^{eq}$, utiliza-se o resolvedor de Newton-Raphson descrito na Seção~\ref{subsec:lbm-equilibrio-g} do Capítulo~\ref{ch:metodo-de-lattice-boltzmann}.

Três decisões distinguem essa solução iterativa de uma implementação convencional de Newton-Raphson, e todas elas decorrem do fato de que o procedimento é executado em cada nó da malha, simultaneamente, e deve ser transposto para processadores de paralelismo massivo.
A primeira é o número fixo de iterações, adotado como $k_{max} = 3$, em lugar de um critério de parada por tolerância.
Um critério de parada faria com que nós vizinhos executassem números distintos de iterações, o que implica divergência entre fluxos de execução e perda de desempenho proporcional ao pior caso do grupo, além de tornar o tempo de execução dependente do estado do escoamento e, por isso, inadequado à comparação entre tecnologias.
A segunda é a partida quente, isto é, o reaproveitamento dos multiplicadores obtidos no passo de tempo anterior como estimativa inicial, o que reduz o número de iterações necessárias, uma vez que a solução varia pouco entre passos consecutivos, ao custo de manter em memória um campo adicional de $D$ componentes por nó.
A terceira é a proteção numérica: o maior expoente é subtraído antes da exponenciação, para evitar transbordamento, e os denominadores $Z$ e $\det \mathbf{J}$ são limitados inferiormente em módulo, de modo que a divisão permaneça finita sem introduzir desvios condicionais no núcleo de cálculo.
Caso o sistema se revele degenerado, adota-se como alternativa a distribuição proporcional aos próprios pesos, $g_i^{eq} = 2\rho E\, W_i(T) / \sum_j W_j(T)$, que satisfaz a restrição de energia total ainda que não a de fluxo, e que é também a expressão empregada na inicialização dos campos.

A etapa de propagação utiliza a técnica de \textit{pull-streaming} combinada com interpolação IDW.
Para otimizar o desempenho, os índices e pesos de interpolação dos quatro vizinhos de cada célula são pré-calculados na fase de inicialização e armazenados em tabelas lineares, eliminando cálculos geométricos redundantes durante a marcha no tempo.

O ajuste de desempenho na versão sequencial foca na eficiência do acesso à memória e na facilitação da vetorização por meio de laços simples e ausência de desvios condicionais complexos dentro do núcleo de computação.
O impacto sobre a complexidade do código é baixo, servindo como uma implementação direta das equações matemáticas do LBM.
O PROGRAMA~\ref{lst:plain-collision} ilustra a estrutura da etapa de colisão, evidenciando o acesso indexado à disposição SoA.

\noindent\begin{minipage}{\linewidth}
\begin{lstlisting}[language=C++, caption={Processamento sequencial da colisão e do equilíbrio.}, label={lst:plain-collision}]
for (Index cell = 0; cell < N; ++cell) {
    const Scalar rho = pRho[cell];
    const Scalar Tv = pT[cell];
    const Scalar ux = pU[0 * N + cell];
    const Scalar uy = pU[1 * N + cell];

    Scalar feq[9], geq[9];
    compute_equilibria(cell, rho, Tv, ux, uy, feq, geq);

    // Aplicacao do operador BGK em cada direcao idc
    for (int idc = 0; idc < 9; ++idc) {
        const Index idx = idc * N + cell;
        pF[idx] += omegaL * (feq[idc] - pF[idx]);
        pG[idx] += omegaL * (geq[idc] - pG[idx]);
    }
}
\end{lstlisting}
\fonte{O autor (2026).}
\end{minipage}

Nesse trecho, $N$ é o número total de células, \texttt{cell} é o índice da célula corrente, \texttt{idc} representa a direção de velocidade discreta, $\omega_L$ é o parâmetro de relaxação e $f_i$ e $g_i$ são as funções de distribuição de massa e de energia, respectivamente.
As tabelas de índices empregadas na propagação já incorporam o tratamento das condições de contorno periódicas ou de parede, de modo que o mapeamento dos vizinhos de cada célula é obtido por leitura direta, sem qualquer desvio condicional no interior do laço.

%---------------------------------------------------------------------------------------

\section{Implementação Utilizando OpenMP}\label{sec:implementacao-computacional-utilizando-openmp}

A implementação em OpenMP contempla tanto o paralelismo em CPUs multinúcleo quanto a aceleração em GPUs por meio de \textit{offloading}.
O modelo de execução adotado para CPUs baseia-se no paradigma \textit{fork-join}, em que o fluxo de execução principal cria uma equipe de fluxos para processar os laços de células da malha.
A paralelização é expressa pela diretiva \texttt{\#pragma omp parallel for}, conforme ilustrado no PROGRAMA~\ref{lst:openmp-cpu}.
Para otimizar o desempenho em CPUs modernas, utiliza-se a cláusula \texttt{simd} para incentivar a vetorização e o escalonamento estático (\texttt{schedule(static)}), que minimiza a latência de gerenciamento ao distribuir a carga de forma equitativa e previsível entre os núcleos.

\noindent\begin{minipage}{\linewidth}
\begin{lstlisting}[language=C++, caption={Diretiva de paralelização OpenMP para CPUs multinúcleo com suporte a SIMD.}, label={lst:openmp-cpu}]
#pragma omp parallel for simd schedule(static)
for (Index cell = 0; cell < N; ++cell) {
    // Calculo da colisao ou da propagacao por celula
    process_cell(cell);
}
\end{lstlisting}
\fonte{O autor (2026).}
\end{minipage}

A gestão de memória na versão de CPU é feita de forma compartilhada no sistema principal, aproveitando a disposição SoA para garantir acessos alinhados.
Na variante voltada a aceleradores, a gestão de memória torna-se explícita.
Utiliza-se a diretiva \texttt{\#pragma omp target data} para mapear os dados na VRAM do dispositivo, enquanto a cláusula \texttt{present} informa ao compilador que os ponteiros já residem no acelerador, evitando transferências redundantes pelo barramento PCIe.
Para maximizar a eficiência em GPUs, a implementação utiliza a cláusula \texttt{collapse(2)} em malhas bidimensionais, ampliando o espaço de iteração disponível para o escalonador do dispositivo.

O mapeamento das etapas do LBM no OpenMP segue a estratégia de fusão de \textit{kernels}, isto é, das funções de cálculo executadas no dispositivo.
Na versão de GPU, as etapas de colisão e de equilíbrio térmico são processadas conjuntamente, assim como a propagação e o cálculo das grandezas macroscópicas.
Essa fusão reduz o tráfego de memória global ao manter variáveis intermediárias em registradores.
Os ajustes de desempenho incluem a fixação do número de fluxos de execução por bloco via \texttt{thread\_limit(64)}, o que otimiza a ocupação dos multiprocessadores ao equilibrar o uso de registradores e a concorrência.

O impacto sobre a complexidade do código é moderado.
Embora a lógica física permaneça próxima da versão sequencial, a introdução de diretivas de \textit{offloading} exige cuidado na gestão do ciclo de vida dos dados e no controle da afinidade dos fluxos de execução (\textit{thread affinity}).
Essa preocupação é particularmente relevante em sistemas de acesso não uniforme à memória (\textit{Non-Uniform Memory Access}, NUMA), arquiteturas em que cada agrupamento de núcleos possui controlador e banco de memória próprios, de modo que o tempo de acesso depende de o dado residir no banco local ou no banco associado a outro agrupamento.
Nesses sistemas, a correta vinculação dos fluxos aos núcleos físicos, aliada à alocação das páginas junto ao núcleo que primeiro as acessa, é fundamental para evitar a degradação de desempenho causada por acessos remotos.

%---------------------------------------------------------------------------------------

\section{Implementação Utilizando OpenACC}\label{sec:implementacao-computacional-utilizando-openacc}

A implementação em OpenACC visa a portabilidade de desempenho com baixo impacto na complexidade do código fonte.
O modelo de execução adotado é descritivo, no qual o programador identifica as regiões de paralelismo e o compilador mapeia as iterações para o dispositivo acelerador.
Utiliza-se a diretiva \texttt{\#pragma acc parallel loop} para distribuir o processamento das células da malha, conforme apresentado no PROGRAMA~\ref{lst:openacc-parallel}.
A estruturação do código em funções marcadas com \texttt{\#pragma acc routine seq} permite a reutilização da lógica de colisão e de equilíbrio sem duplicação de código, preservando a modularidade entre as versões de CPU e de GPU.

\noindent\begin{minipage}{\linewidth}
\begin{lstlisting}[language=C++, caption={Diretivas OpenACC para paralelização em GPU.}, label={lst:openacc-parallel}]
#pragma acc parallel loop present(pF, pG, pRho, pT, pU) \
    num_gangs((N + 63) / 64) vector_length(64)
for (int cell = 0; cell < N; ++cell) {
    // Processamento paralelo da celula no dispositivo
    collision_kernel(cell);
}
\end{lstlisting}
\fonte{O autor (2026).}
\end{minipage}

A gestão de memória no OpenACC é feita de forma integrada ao ecossistema do NVIDIA HPC SDK.
Embora o OpenACC possua diretivas próprias para gestão de dados (\texttt{copyin}, \texttt{copyout}), a implementação utiliza a interface de alocação de memória de dispositivo do OpenMP (\texttt{omp\_target\_alloc}) para reservar regiões persistentes na GPU.
Essa decisão técnica permite uma interoperabilidade eficiente entre os modelos de paralelização.
O uso da cláusula \texttt{deviceptr} informa ao compilador que os dados já residem na VRAM, eliminando o rastreamento de presença e reduzindo a latência de lançamento dos \textit{kernels}.

O mapeamento das etapas do LBM no OpenACC segue a técnica de fusão de laços para minimizar o tráfego de dados.
A etapa de propagação utiliza uma abordagem de interpolação IDW calculada em tempo de execução (\textit{on-the-fly}), o que evita o armazenamento de grandes tabelas de índices que poderiam exceder a capacidade da VRAM em malhas extensas.
A lógica de interpolação é tornada livre de desvios (\textit{branchless}) por meio do uso de funções transcendentais e sentinelas numéricas, garantindo que todos os fluxos de execução de um grupo processem a mesma sequência de instruções.

Os ajustes de desempenho focam na ocupação dos multiprocessadores e na eficiência dos acessos à memória.
Utiliza-se a cláusula \texttt{vector\_length(64)} para alinhar o tamanho do grupo de execução com as características do \textit{hardware} alvo, evitando o transbordamento de registradores (\textit{register spilling}) para a memória global.
O impacto sobre a complexidade do código é mínimo, pois o OpenACC permite manter uma única base de código C++ para diferentes arquiteturas, delegando ao compilador as tarefas de especialização e de otimização de baixo nível.

%---------------------------------------------------------------------------------------

\section{Implementação Utilizando CUDA}\label{sec:implementacao-computacional-utilizando-cuda}

A implementação em CUDA representa a versão de maior desempenho para aceleradores NVIDIA, utilizando um modelo de programação de baixo nível que permite controle fino sobre o \textit{hardware}.
O modelo de execução adotado é o SIMT, no qual a malha computacional é mapeada em uma hierarquia de blocos e fluxos de execução.
Cada célula da malha é processada por um fluxo de execução individual, agrupados em blocos de tamanho fixo (64 fluxos por bloco), o que garante alta ocupação dos multiprocessadores (\textit{Streaming Multiprocessors}, SMs).

A gestão de memória é feita de forma explícita na VRAM da GPU utilizando a API \texttt{cudaMalloc}.
Para otimizar o acesso aos dados, utiliza-se a disposição SoA, permitindo o coalescimento de memória (\textit{memory coalescing}) em que as requisições de múltiplos fluxos de um \textit{warp} são atendidas em uma única transação de barramento.
O \textit{warp} é o grupo de trinta e dois fluxos de execução que o multiprocessador escalona como unidade indivisível, executando a mesma instrução sobre dados distintos; quando fluxos de um mesmo \textit{warp} seguem caminhos diferentes em um desvio condicional, os caminhos são percorridos em sequência, situação designada divergência de \textit{warp}, que reduz o aproveitamento das unidades de execução.
As tabelas de velocidades da rede D2Q9 são armazenadas em memória constante, qualificada por \texttt{\_\_constant\_\_}, o que reduz a latência de acesso para dados lidos simultaneamente por todos os fluxos de execução.
Adicionalmente, utiliza-se a técnica de duplo armazenamento temporário (\textit{double buffering}) para as distribuições, eliminando dependências de dados e condições de corrida durante a etapa de propagação.

O mapeamento das etapas do LBM utiliza a fusão de \textit{kernels} para reduzir a latência de lançamento e o tráfego com a memória global.
O PROGRAMA~\ref{lst:cuda-parallel} ilustra o \textit{kernel} de colisão, destacando o uso de registradores e o desdobramento de laços por meio de \texttt{\#pragma unroll}.

\noindent\begin{minipage}{\linewidth}
\begin{lstlisting}[language=C++, caption={\textit{Kernel} CUDA para a etapa de colisão.}, label={lst:cuda-parallel}]
template<typename Scalar>
__global__ void collisionKernel(Scalar* pF, Scalar* pG, int N) {
    int cell = blockIdx.x * blockDim.x + threadIdx.x;
    if (cell < N) {
        Scalar feq[9], geq[9];
        // Calculo de feq e geq em registradores
        compute_equilibria(cell, feq, geq);

        #pragma unroll
        for (int idc = 0; idc < 9; ++idc) {
            int idx = idc * N + cell;
            pF[idx] += omegaL * (feq[idc] - pF[idx]);
        }
    }
}
\end{lstlisting}
\fonte{O autor (2026).}
\end{minipage}

Os ajustes de desempenho incluem a eliminação de divergência de \textit{warp} por meio de lógica livre de desvios na solução de Newton-Raphson e na interpolação IDW.
A interpolação utiliza a aproximação \texttt{exp(-p*log(r2+kTiny))} para evitar testes condicionais nas fronteiras da malha.
O impacto sobre a complexidade do código é elevado, pois exige a manutenção de arquivos \texttt{.cu} e o uso do compilador \texttt{nvcc}, além da gestão manual da hierarquia de memória e do alinhamento de dados.

%---------------------------------------------------------------------------------------

\section{Implementação Utilizando OpenCL}\label{sec:implementacao-computacional-utilizando-opencl}

O módulo OpenCL provê portabilidade funcional para diversos aceleradores, incluindo GPUs de diferentes fabricantes e CPUs por meio do \textit{runtime} POCL (\textit{Portable Computing Language}).
O modelo de execução baseia-se em uma plataforma que gerencia dispositivos aceleradores, nos quais os \textit{kernels} são despachados para execução em um espaço de índices global (\texttt{NDRange}).
Diferentemente do CUDA, o OpenCL utiliza compilação em tempo de execução (JIT), em que o código fonte dos \textit{kernels} é lido de arquivos externos e compilado no momento da inicialização do miniaplicativo.

A gestão de memória é feita por meio de objetos de \textit{buffer} (\texttt{cl::Buffer}) que representam regiões de memória no dispositivo.
A transferência de dados entre o sistema principal e o acelerador ocorre de forma explícita via filas de comandos (\texttt{cl::CommandQueue}).
Um recurso avançado empregado nesta versão é a extensão de \textit{buffers} de comandos, identificada por \texttt{cl\_khr\_command\_buffer}, que permite registrar a sequência de chamadas de \textit{kernels} da marcha no tempo em uma estrutura persistente.
Essa técnica reduz drasticamente a carga sobre o processador principal ao evitar o reenvio individual de comandos a cada passo de tempo, conforme ilustrado no PROGRAMA~\ref{lst:opencl-cpp}.

\noindent\begin{minipage}{\linewidth}
\begin{lstlisting}[language=C++, caption={Uso de \textit{buffers} de comandos para redução de latência no OpenCL.}, label={lst:opencl-cpp}]
// Registro dos comandos da marcha no tempo
pfn_ndrange(cbEven, queue, NULL, collisionKernel, 1, NULL, &N, &LWS, 0, NULL, NULL, NULL);
pfn_ndrange(cbEven, queue, NULL, streamKernel, 1, NULL, &N, &LWS, 0, NULL, NULL, NULL);
pfn_finalize(cbEven);

// Execucao eficiente por passo de tempo
pfn_enqueue(0, NULL, cbEven, 0, NULL, NULL);
\end{lstlisting}
\fonte{O autor (2026).}
\end{minipage}

O mapeamento das etapas do LBM segue a mesma lógica de disposição SoA e de fusão de \textit{kernels} adotada no CUDA.
A linguagem OpenCL C é utilizada para expressar a física do modelo, permitindo otimizações específicas como o uso de tipos vetoriais e extensões de precisão dupla.
Os ajustes de desempenho incluem a parametrização do tamanho do grupo de trabalho (\textit{work-group size}) e a sincronização fina entre \textit{kernels} por meio de eventos (\texttt{cl::Event}).

O impacto sobre a complexidade do código é muito elevado.
A necessidade de gerenciar contextos, filas, compilação JIT e a verbosidade intrínseca da API OpenCL tornam a manutenção mais desafiadora.
Contudo, essa tecnologia garante que o miniaplicativo possa ser executado em praticamente qualquer acelerador moderno, fornecendo a base necessária para a portabilidade entre diferentes fabricantes de \textit{hardware}.

%---------------------------------------------------------------------------------------

\section{Implementação Utilizando Kokkos}\label{sec:implementacao-computacional-utilizando-kokkos}

A implementação em Kokkos baseia-se em um modelo de programação de fonte única (\textit{single-source}), permitindo que o mesmo código C++ seja executado em diversas arquiteturas apenas alterando o espaço de execução em tempo de compilação.
O modelo de execução adotado utiliza políticas de intervalo (\texttt{RangePolicy}) e o despacho paralelo de laços por meio de \texttt{Kokkos::parallel\_for}, conforme ilustrado no PROGRAMA~\ref{lst:kokkos-parallel}.
A classe resolvedora é parametrizada pelo espaço de execução, o que possibilita a geração automática de versões para CPUs multinúcleo e para GPUs NVIDIA, AMD ou Intel.

\noindent\begin{minipage}{\linewidth}
\begin{lstlisting}[language=C++, caption={Despacho paralelo sobre Views em Kokkos.}, label={lst:kokkos-parallel}]
using ScalarView = Kokkos::View<Scalar*, typename ExecSpace::memory_space>;

Kokkos::parallel_for(
    Kokkos::RangePolicy<ExecSpace>(0, N),
    KOKKOS_LAMBDA(const Index cell) {
        // Logica de colisao e de equilibrio fundida
        compute_collision(cell);
    });
\end{lstlisting}
\fonte{O autor (2026).}
\end{minipage}

A gestão de memória é feita por meio de \textit{Views}, abstrações de tensores multidimensionais que gerenciam a alocação e o ciclo de vida dos dados nos diferentes espaços de memória.
O Kokkos utiliza contagem de referências para a liberação automática de recursos, simplificando a manutenção em relação ao CUDA.
Para o acesso aos dados, utiliza-se o espelhamento (\textit{mirroring}), em que cópias sincronizadas dos dados são mantidas no sistema principal para facilitar a exportação de resultados, enquanto a marcha no tempo ocorre exclusivamente na memória do dispositivo acelerador.

O mapeamento das etapas do LBM no Kokkos segue a estratégia de núcleos fundidos, otimizando a localidade de dados e o uso de registradores.
Os ajustes de desempenho incluem a especialização de políticas de execução para cada alvo e a garantia de acessos coalescidos por meio da configuração da disposição de memória das \textit{Views}.
O impacto sobre a complexidade do código é reduzido, pois a unificação do código de paralelismo elimina a necessidade de múltiplas implementações para diferentes aceleradores, facilitando a inclusão de novos modelos físicos no miniaplicativo.

%---------------------------------------------------------------------------------------

\section{Implementação Utilizando RAJA}\label{sec:implementacao-computacional-utilizando-raja}

A integração com o RAJA permite escrever laços de repetição agnósticos à plataforma por meio de uma camada de abstração de execução baseada em políticas.
Diferentemente do Kokkos, o RAJA foca exclusivamente no encapsulamento dos laços paralelos, não fornecendo estruturas de dados próprias.
O modelo de execução adotado utiliza a função \texttt{RAJA::forall}, que recebe uma política de execução e uma função \textit{lambda} contendo a lógica da célula, conforme demonstrado no PROGRAMA~\ref{lst:raja-backend}.
Essa abordagem permite compor diferentes estratégias de execução, como sequencial, OpenMP ou CUDA, mantendo a integridade da lógica física.

\noindent\begin{minipage}{\linewidth}
\begin{lstlisting}[language=C++, caption={Definição de backend e uso de RAJA::forall.}, label={lst:raja-backend}]
struct DeviceBackend {
    using ExecPolicy = RAJA::cuda_exec<64>;
    using Memory = DeviceMemory;
};

RAJA::forall<typename Backend::ExecPolicy>(
    RAJA::RangeSegment(0, N),
    [=] RAJA_HOST_DEVICE (Index cell) {
        // Logica de colisao portavel
        compute_collision(cell);
    });
\end{lstlisting}
\fonte{O autor (2026).}
\end{minipage}

A gestão de memória na implementação RAJA é delegada à classe \texttt{DeviceBuffer}, que utiliza características de memória (\textit{memory traits}) para selecionar o local de alocação (RAM ou VRAM) com base no \textit{backend} escolhido.
Essa separação entre a execução e os dados exige que o programador gerencie explicitamente o local onde os ponteiros residem, o que aumenta a flexibilidade, mas também a responsabilidade de manter a coerência entre o sistema principal e o acelerador.

O mapeamento das etapas do LBM no RAJA segue a estrutura modular das versões anteriores, com núcleos fundidos para maximizar o desempenho.
Os ajustes de desempenho são feitos pela seleção criteriosa das políticas de execução e pelo ajuste dos tamanhos de bloco para GPUs.
O impacto sobre a complexidade do código é reduzido na camada de lógica física, mas ligeiramente superior na camada de infraestrutura devido à necessidade de gerenciar explicitamente os \textit{backends} de memória e de execução.
A flexibilidade do RAJA permite que o miniaplicativo se adapte rapidamente a novas arquiteturas de HPC sem a necessidade de reescrever o núcleo computacional.

%---------------------------------------------------------------------------------------

\section{Implementação Utilizando MPI}\label{sec:implementacao-computacional-utilizando-mpi}

Para viabilizar simulações de larga escala, que excedem a capacidade de memória e de processamento de um único nó de computação, utiliza-se o padrão MPI para paralelismo em memória distribuída.
O modelo de execução adotado baseia-se na decomposição de domínio unidimensional ao longo do eixo longitudinal ($x$), técnica particularmente eficiente para o estudo de tubos de choque, nos quais os gradientes mais intensos ocorrem nessa direção.
Cada processo MPI é responsável por uma fatia local da malha, executando a lógica do LBM de forma independente na maior parte do domínio.

A gestão de memória no ambiente distribuído fundamenta-se no uso de camadas de células fantasma (\textit{ghost cells}).
Essas camadas são regiões de memória adicionais nas fronteiras de cada subdomínio que armazenam cópias das distribuições pertencentes aos processos vizinhos imediatos.
Para a exportação de resultados globais, utiliza-se a função \texttt{MPI\_Allgatherv}, que reconstrói os campos macroscópicos completos a partir das fatias locais processadas por cada processo.
A \autoref{fig:mpi-decomposicao} ilustra essa decomposição e as zonas de troca entre subdomínios.

\begin{figure}[H]
    \centering
    \caption{Decomposição unidimensional do domínio entre processos MPI, com as camadas de células fantasma utilizadas na troca de distribuições entre subdomínios vizinhos.}
    \label{fig:mpi-decomposicao}
    \fbox{\begin{minipage}{0.85\textwidth}
        \centering
        \vspace{2cm}
        [Espaço reservado para figura: decomposição unidimensional do domínio do tubo de choque ao longo do eixo longitudinal, em quatro subdomínios atribuídos a processos MPI distintos.
        Nas fronteiras entre subdomínios devem ser destacadas as camadas de células fantasma e as setas de envio e de recebimento das distribuições entre processos vizinhos.]
        \vspace{2cm}
    \end{minipage}}
    \fonte{O autor (2026).}
\end{figure}

O mapeamento das etapas do LBM no MPI é projetado para maximizar a sobreposição entre comunicação e cálculo.
Após o processamento das células de fronteira, cada processo inicia a troca das distribuições por meio de comunicações não bloqueantes (\texttt{MPI\_Isend} e \texttt{MPI\_Irecv}).
Enquanto os dados trafegam pela rede, o processo continua o cálculo da colisão nas células internas do domínio, cujos dados não dependem de informações externas.
Essa estratégia, exemplificada no PROGRAMA~\ref{lst:mpi-pack}, minimiza o tempo de espera ociosa e melhora a escalabilidade forte do miniaplicativo.

\noindent\begin{minipage}{\linewidth}
\begin{lstlisting}[language=C++, caption={Comunicação não bloqueante e sobreposição de cálculo no MPI.}, label={lst:mpi-pack}]
// Inicio da troca de fronteiras (nao bloqueante)
startExchangeGhosts(fCur, gCur);

// Calculo da colisao no miolo do dominio (sobreposicao)
collisionInnerDomain(fCur, gCur);

// Finalizacao da comunicacao
finishExchangeGhosts();
\end{lstlisting}
\fonte{O autor (2026).}
\end{minipage}

Os ajustes de desempenho incluem o empacotamento explícito dos dados de fronteira em regiões contíguas de memória, o que é necessário devido à disposição SoA.
Adicionalmente, utilizam-se tipos de dados derivados (\texttt{MPI\_Type\_vector}) para otimizar a transferência de colunas de dados não contíguos em malhas mais complexas.
O impacto sobre a complexidade do código é elevado, pois exige a gestão de subdomínios, a sincronização explícita e o tratamento de condições de contorno em ambientes distribuídos.
Contudo, essa implementação é essencial para atingir o nível de desempenho exigido em sistemas de HPC de grande porte.

%---------------------------------------------------------------------------------------

\section{Compilação, Configuração e Execução}\label{sec:compilacao-configuracao}

O processo de construção do miniaplicativo utiliza o sistema CMake para automatizar a detecção de dependências e a configuração do ambiente de compilação.
O CMake identifica a presença de compiladores C++ compatíveis com o padrão C++17 e busca bibliotecas externas como Eigen, yaml-cpp e implementações de MPI e OpenCL.
Para tecnologias que exigem compiladores específicos, como o \texttt{nvcc} para CUDA ou o NVIDIA HPC SDK para OpenACC e OpenMP-target, o sistema verifica a disponibilidade das ferramentas e habilita condicionalmente os respectivos \textit{backends}.
A biblioteca Eigen e o leitor de arquivos YAML são obtidos automaticamente via \texttt{FetchContent}, garantindo que as dependências estejam na versão correta.

A configuração da simulação é centralizada no arquivo \texttt{config.yaml}, que permite ao usuário definir os parâmetros físicos e numéricos sem a necessidade de recompilar o código.
O PROGRAMA~\ref{lst:config-yaml} apresenta o conteúdo padrão do arquivo de configuração, estruturado em quatro grupos principais: \texttt{Fluid}, \texttt{Mesh}, \texttt{Control} e \texttt{Performance}.

\noindent\begin{minipage}{\linewidth}
\begin{lstlisting}[language=bash, caption={Arquivo de configuração config.yaml do miniaplicativo.}, label={lst:config-yaml}]
Fluid:
  densityL: 1.0
  pressureL: 1.0
  densityR: 0.125
  pressureR: 0.1
  viscosity: 1.0e-2
  prandtl: 0.71
  gamma: 1.4
Mesh:
  lx: 1.0
  ly: 1.0
  nx: 50000
  ny: 500
Control:
  tmax: 0.2
  Ux: 0.25
  Uy: 0.0
  printStep: 20
Performance:
  target: CPU
  backend: Plain
  cores: 0
\end{lstlisting}
\fonte{O autor (2026).}
\end{minipage}

O grupo \texttt{Fluid} define as propriedades termodinâmicas e de transporte do fluido.
Os parâmetros \texttt{densityL}, \texttt{pressureL}, \texttt{densityR} e \texttt{pressureR} estabelecem os estados inicial à esquerda e à direita da ruptura do diafragma no problema de Sod, sendo fundamentais para a caracterização do choque e do leque de rarefação.
A constante \texttt{gamma} representa a razão de calores específicos, enquanto \texttt{viscosity} e \texttt{prandtl} governam a dissipação de momento e de calor, respectivamente.

O grupo \texttt{Mesh} especifica as dimensões físicas do domínio (\texttt{lx}, \texttt{ly}) e a resolução da malha (\texttt{nx}, \texttt{ny}).
Como o problema de Sod possui gradientes apenas na direção longitudinal, a resolução transversal não altera a solução física, mas determina o comprimento das sequências contíguas percorridas em cada linha da malha, razão pela qual é mantida na proporção $ny = nx/100$ adotada na Seção~\ref{subsec:solucao-numerica-utilizando-o-metodo-lbm}.
O grupo \texttt{Control} define o tempo máximo de simulação (\texttt{tmax}), as componentes da velocidade de referência da rede deslocada (\texttt{Ux}, \texttt{Uy}) e a frequência de gravação dos resultados (\texttt{printStep}).

Por fim, o grupo \texttt{Performance} determina o dispositivo de execução (\texttt{target}), que aceita os valores \texttt{CPU} ou \texttt{GPU}, e a tecnologia de paralelização (\texttt{backend}), que aceita os valores \texttt{Plain}, \texttt{OpenMP}, \texttt{OpenACC}, \texttt{CUDA}, \texttt{OpenCL}, \texttt{Kokkos}, \texttt{RAJA} ou \texttt{MPI}.
O parâmetro \texttt{cores} permite definir o número de fluxos de execução em CPUs, sendo o valor 0 utilizado para detecção automática com base no \textit{hardware} disponível.
